\documentclass{article}
\usepackage{graphicx} % Required for inserting images
\usepackage{amssymb} % For \mathbb
\usepackage{amsmath}

\begin{document}

\section{Efficient Computation of Greeks via Fast Fourier Transform}
\subsection{Analytic Derivation of the Delta via FFT}

Following the framework of Carr and Madan (1999)[cite: 178], the value of a European call option $C_T(k)$ at log-strike $k$ is obtained by inverting the Fourier transform of the damped price $c_T(k) = e^{\alpha k}C_T(k)$. The pricing formula is given by:

\begin{equation}
    C_T(k) = \frac{e^{-\alpha k}}{\pi} \int_{0}^{\infty} e^{-ivk} \psi_T(v) \, dv
    \label{eq:carr_madan_inversion}
\end{equation}

where $\psi_T(v)$, the Fourier transform of $c_T(k)$, is defined analytically in terms of the characteristic function $\phi_T(u) = \mathbb{E}[e^{iu s_T}]$ as:

\begin{equation}
    \psi_T(v) = \frac{e^{-rT} \phi_T(v - (\alpha+1)i)}{\alpha^2 + \alpha - v^2 + i(2\alpha+1)v}
    \label{eq:psi_def}
\end{equation}

To derive the option Delta, $\Delta = \frac{\partial C_T}{\partial S_0}$, we exploit the separability of the characteristic function with respect to the initial spot price $S_0$. Letting $X_T$ denote the process increments such that $s_T = \ln S_0 + X_T$, the characteristic function factorizes as:

\begin{equation}
    \phi_T(u) = \mathbb{E}[e^{i u s_T}] = \mathbb{E}[e^{i u (\ln(S_0) + X_T)}] = e^{i u \ln(S_0)} \mathbb{E}[e^{i u X_T}] = S_0^{i u} \phi_T^{X}(u)
\end{equation}

Differentiating with respect to $S_0$ yields:

\begin{equation}
    \frac{\partial \phi_T(u)}{\partial S_0} = i u S_0^{i u - 1} \phi_T^{X}(u) = \frac{i u}{S_0} \phi_T(u)
    \label{eq:cf_deriv_delta}
\end{equation}

Since the integration in Equation 
\eqref{eq:carr_madan_inversion} 
is linear and the integral converges, we can differentiate under the integral sign. The partial derivative of the integrand $\psi_T(v)$ with respect to $S_0$ is obtained by applying the chain rule to Equation \eqref{eq:psi_def}.
Let $u = v - (\alpha+1)i$ be the argument of the characteristic function in the pricing formula. Substituting the derivative from Equation \eqref{eq:cf_deriv_delta}:

\begin{equation}
    \frac{\partial \psi_T(v)}{\partial S_0} = \frac{e^{-rT}}{\alpha^2 + \alpha - v^2 + i(2\alpha+1)v} \cdot \frac{\partial \phi_T(v - (\alpha+1)i)}{\partial S_0}
\end{equation}

\begin{equation}
    \frac{\partial \psi_T(v)}{\partial S_0} = \psi_T(v) \cdot \frac{i(v - (\alpha+1)i)}{S_0}
\end{equation}

Consequently, $\psi_{\Delta}(v)$ is given by a linear modulation of the pricing transform:

\begin{equation}
    \psi_{\Delta}(v) = \frac{\partial \psi_T(v)}{\partial S_0} = \frac{\alpha + 1 + iv}{S_0} \psi_T(v)
\end{equation}

The Delta is efficiently computed for all strikes $k$ via a single Fast Fourier Transform:

\begin{equation}
    \Delta(k) = \frac{e^{-\alpha k}}{\pi} \int_{0}^{\infty} e^{-ivk} \psi_{\Delta}(v) \, dv
\end{equation}



\subsection{Analytic Derivation of the Gamma}

The option Gamma, defined as $\Gamma = \frac{\partial^2 C_T}{\partial S_0^2}$, is derived analogously by differentiating the characteristic function a second time with respect to the spot price. Differentiating \eqref{eq:cf_deriv_delta} with respect to $S_0$ yields:

\begin{equation}
    \frac{\partial^2 \phi_T(u)}{\partial S_0^2} = \frac{\partial}{\partial S_0} \left( \frac{iu}{S_0} \phi_T(u) \right)
\end{equation}

Applying the product rule:

\begin{equation}
    \frac{\partial^2 \phi_T(u)}{\partial S_0^2} = -\frac{iu}{S_0^2}\phi_T(u) + \frac{iu}{S_0} \left( \frac{iu}{S_0} \phi_T(u) \right)
\end{equation}

\begin{equation}
    \frac{\partial^2 \phi_T(u)}{\partial S_0^2} = \frac{iu(iu - 1)}{S_0^2} \phi_T(u) = -\frac{u^2 + iu}{S_0^2} \phi_T(u)
\end{equation}

To obtain the Fourier transform of the Gamma, $\psi_{\Gamma}(v)$, we apply this weighting factor to the pricing integrand $\psi_T(v)$. Recalling that the argument of the characteristic function in the Carr-Madan formula is $u = v - (\alpha+1)i$, the term $iu$ becomes:
$$ i(v - i\alpha - i) = iv + \alpha + 1 $$

Substituting this into the second derivative expression, the integrand for the Gamma becomes:

\begin{equation}
    \psi_{\Gamma}(v) = \frac{1}{S_0^2} \left[ (iv + \alpha + 1)(iv + \alpha) \right] \psi_T(v)
\end{equation}

The Gamma values are subsequently recovered via the inverse FFT:

\begin{equation}
    \Gamma(k) = \frac{e^{-\alpha k}}{\pi} \int_{0}^{\infty} e^{-ivk} \psi_{\Gamma}(v) \, dv
\end{equation}
\subsection{Analytic Derivation of the Vega (Black-Scholes Case)}

The Vega, $\mathcal{V} = \frac{\partial C_T}{\partial \sigma}$, is model-dependent as it requires differentiating the characteristic function with respect to the volatility parameter $\sigma$. Assuming the underlying asset follows a Geometric Brownian Motion (Black-Scholes model), the characteristic function of the log-price is given by:

\begin{equation}
    \phi_T^{BS}(u) = \exp \left( i u \left(\ln S_0 + (r - q - \frac{1}{2}\sigma^2)T \right) - \frac{1}{2} \sigma^2 u^2 T \right)
\end{equation}

Differentiating $\phi_T^{BS}(u)$ with respect to $\sigma$, we obtain:

\begin{equation}
    \frac{\partial \phi_T^{BS}(u)}{\partial \sigma} = \phi_T^{BS}(u) \cdot \frac{\partial}{\partial \sigma} \left[ i u (r - q - \frac{\sigma^2}{2})T - \frac{1}{2} \sigma^2 u^2 T \right]
\end{equation}

\begin{equation}
    \frac{\partial \phi_T^{BS}(u)}{\partial \sigma} = \phi_T^{BS}(u) \cdot \left[ -iu\sigma T - \sigma u^2 T \right]
\end{equation}

\begin{equation}
    \frac{\partial \phi_T^{BS}(u)}{\partial \sigma} = -\sigma T (u^2 + iu) \phi_T^{BS}(u)
\end{equation}

Thus, the Fourier transform of the Vega, $\psi_{\mathcal{V}}(v)$, is obtained by modulating the pricing integrand with the factor $-\sigma T (u^2 + iu)$, evaluated at $u = v - (\alpha+1)i$:

\begin{equation}
    \psi_{\mathcal{V}}(v) = -\sigma T \left[ (v - (\alpha+1)i)^2 + i(v - (\alpha+1)i) \right] \psi_T(v)
\end{equation}

Using the identity $u^2 + iu = iu(iu-1)$ derived for the Gamma, this can also be expressed as:

\begin{equation}
    \psi_{\mathcal{V}}(v) = -\sigma T \left[ (iv + \alpha + 1)(iv + \alpha) \right] \psi_T(v)
\end{equation}

Finally, the Vega is computed via the inverse FFT:

\begin{equation}
    \mathcal{V}(k) = \frac{e^{-\alpha k}}{\pi} \int_{0}^{\infty} e^{-ivk} \psi_{\mathcal{V}}(v) \, dv
\end{equation}


\subsection{Computational Efficiency}

The FFT-based evaluation of Greeks offers significant computational advantages over traditional finite difference methods or direct integration. We highlight two primary benefits:

\begin{itemize}
    \item \textbf{Algorithmic Complexity:} The Fast Fourier Transform reduces the computational cost to $O(N \log_2 N)$, compared to $O(N^2)$ for direct quadrature. This allows for the \textit{simultaneous} valuation of sensitivities across a dense grid of $N$ strikes in a single execution, rendering the marginal cost of pricing additional strikes negligible.

    \item \textbf{Numerical Precision:} The derivation of analytic expressions for $\psi_{\Delta}$, $\psi_{\Gamma}$, and $\psi_{\mathcal{V}}$ in the frequency domain eliminates the need for numerical differentiation (finite differences). This approach avoids the computational overhead of multiple pricing passes and circumvents the truncation and round-off errors inherent in approximating derivatives via perturbation.
\end{itemize}

\end{document}